\documentclass[aspectratio=169,11pt]{beamer}

\usetheme{Madrid}
\usecolortheme{dolphin}
\setbeamertemplate{navigation symbols}{}
\setbeamertemplate{itemize item}[circle]
\setbeamertemplate{itemize subitem}[triangle]

\usepackage{amsmath}
\usepackage{graphicx}

\title[Beamer with LaTeX]{Beamer: Creating Presentations with LaTeX}
\subtitle{A Short Introduction}
\author{Ridhima Garg}
\institute{2025UCP1648}
\date{Technical Writing }

\begin{document}

% Slide 1: Title
\begin{frame}
    \titlepage
\end{frame}

% Slide 2: Introduction
\begin{frame}{What is Beamer?}
    \begin{block}{Definition}
        Beamer is a LaTeX document class used to create structured, high-quality PDF presentations.
    \end{block}

    \vspace{0.25cm}
    \begin{itemize}
        \item Slides are created using the \texttt{frame} environment.
        \item Content is written as LaTeX source code.
        \item Mathematical expressions, tables, and figures can be included.
        \item The final presentation is compiled into a PDF.
    \end{itemize}
\end{frame}

% Slide 3: Core building blocks
\begin{frame}{The Building Blocks of a Beamer Presentation}
    \begin{columns}[T,onlytextwidth]
        \column{0.53\textwidth}
        \begin{block}{Document setup}
            \small
            \texttt{\textbackslash documentclass\{beamer\}}\\[0.12cm]
            \texttt{\textbackslash usetheme\{Madrid\}}\\[0.12cm]
            \texttt{\textbackslash title\{...\}}\\
            \texttt{\textbackslash author\{...\}}
        \end{block}
        \begin{block}{Slide content}
            \small
            \texttt{\textbackslash begin\{frame\}}\\
            \texttt{\textbackslash end\{frame\}}
        \end{block}

        \column{0.42\textwidth}
        \begin{exampleblock}{Example equation}
            \[
                E = mc^2
            \]
            \centering
            \small Mathematical notation works naturally in Beamer.
        \end{exampleblock}
        \begin{alertblock}{Remember}
            A frame usually represents one slide.
        \end{alertblock}
    \end{columns}
\end{frame}

% Slide 4: Workflow and conclusion
\begin{frame}{From Source Code to Slides}
    \begin{enumerate}
        \item Choose the Beamer document class and a theme.
        \item Add presentation details such as title and author.
        \item Create frames and organize content with lists or blocks.
        \item Compile the \texttt{.tex} file to generate a PDF.
    \end{enumerate}

    \vspace{0.35cm}
    \begin{block}{Key takeaway}
        Beamer combines the precision of LaTeX with a structured way to design academic presentations.
    \end{block}

    \vfill
    \centering
    \Large\textbf{Thank You!}
\end{frame}

\end{document}